% GENERATED FILE. Edit canonical lessons, then run npm run generate:books.
% canonical-source-hash: fnv1a64:d29767cbc63268d4
% canonical-lessons: TA-C335-murrukaiyitu, TA-C335-akali, TA-C335-curankam, TA-C335-aran, TA-C335-orran

\chapter{A Siege}
\label{ch:ta-a-siege}
\hlchaptermodality{335}

\begin{chapteropening}
\textbf{By the end of this chapter:} \emph{I can describe a siege and how a fort was defended.}

Complete the last lesson of chapter 335: I can describe a siege and how a fort was defended.
\end{chapteropening}

\section[\texorpdfstring{\ta{முற்றுகையிடு} (muṟṟukaiyiṭu)}{முற்றுகையிடு (muṟṟukaiyiṭu)}]{\ta{முற்றுகையிடு} (muṟṟukaiyiṭu) — to besiege}

\label{lesson:TA-C335-murrukaiyitu}



\begin{quote}
Before the new one: say the Tamil for a war drum, then the Tamil for a commander.
\end{quote}

\subsection*{You'll want to know: \ta{முற்றுகையிடு}}
\textbf{\ta{முற்றுகையிடு}} — \emph{muṟṟukaiyiṭu} — \textquotedblleft{}to besiege\textquotedblright{}.

\begin{quote}
\ta{படை} \ta{கோட்டையை} \ta{ஒரு} \ta{மாதம்} \ta{முற்றுகையிட்டது}.
\end{quote}

\begin{quote}
\emph{paṭai kōṭṭaiyai oru mātam muṟṟukaiyiṭṭatu.}
\end{quote}

\begin{quote}
The army besieged the fort for a month.
\end{quote}

Past (I …): \ta{முற்றுகையிட்டேன்}, \emph{muṟṟukaiyiṭṭēṉ}.

\ta{முற்றுகை}, a siege + \ta{இடு}, to put.

\begin{cousinweb}
\textbf{\ta{முற்றுகையிடு}} (\emph{muṟṟukaiyiṭu}) is built on native \ta{முற்றுகை}, \textquotedblleft{}blockade, siege\textquotedblright{}, from \ta{முற்று}, \textquotedblleft{}to surround, besiege\textquotedblright{} (DEDR 5018), with \ta{இடு}, \textquotedblleft{}to put\textquotedblright{}.
\end{cousinweb}

\begin{grammarlens}[title={What you've built}]
One more word for telling what happened.
\end{grammarlens}

\subsection*{Your turn}
\begin{itemize}
\raggedright
  \item \emph{Say it:} \emph{muṟṟukaiyiṭu}
  \item \emph{Say it:} \emph{muṟṟukaiyiṭu}, once more
  \item \emph{Say it:} \emph{paṭai kōṭṭaiyai oru mātam muṟṟukaiyiṭṭatu.} — The army besieged the fort for a month.
  \item \emph{Recall:} say the Tamil for a war drum, then the Tamil for a commander, then say \emph{muṟṟukaiyiṭu} again
\end{itemize}

\begin{tcolorbox}[breakable,colback=teal!4,colframe=teal!35!black,title={Before you move on}]
What does \ta{முற்றுகையிடு} mean? (\textquotedblleft{}To besiege\textquotedblright{}.) Say it once more.
\end{tcolorbox}
\section[\texorpdfstring{\ta{அகழி} (akaḻi)}{அகழி (akaḻi)}]{\ta{அகழி} (akaḻi) — a moat}

\label{lesson:TA-C335-akali}



\begin{quote}
Before the new one: say the Tamil for a commander, then the Tamil for to besiege.
\end{quote}

\subsection*{You'll want to know: \ta{அகழி}}
\textbf{\ta{அகழி}} — \emph{akaḻi} — \textquotedblleft{}a moat\textquotedblright{}.

\begin{quote}
\ta{கோட்டையைச்} \ta{சுற்றி} \ta{ஆழமான} \ta{அகழி} \ta{இருந்தது}.
\end{quote}

\begin{quote}
\emph{kōṭṭaiyaic cuṟṟi āḻamāṉa akaḻi iruntatu.}
\end{quote}

\begin{quote}
There was a deep moat around the fort.
\end{quote}

From \ta{அகழ்}, to dig.

\begin{cousinweb}
\textbf{\ta{அகழி}} (\emph{akaḻi}) is native Dravidian: DEDR 11 gives it under \ta{அகழ்}, \textquotedblleft{}to excavate, dig out\textquotedblright{}.
\end{cousinweb}

\begin{grammarlens}[title={What you've built}]
One more word for telling what happened.
\end{grammarlens}

\subsection*{Your turn}
\begin{itemize}
\raggedright
  \item \emph{Say it:} \emph{akaḻi}
  \item \emph{Say it:} \emph{akaḻi}, once more
  \item \emph{Say it:} \emph{kōṭṭaiyaic cuṟṟi āḻamāṉa akaḻi iruntatu.} — There was a deep moat around the fort.
  \item \emph{Recall:} say the Tamil for a commander, then the Tamil for to besiege, then say \emph{akaḻi} again
\end{itemize}

\begin{tcolorbox}[breakable,colback=teal!4,colframe=teal!35!black,title={Before you move on}]
What does \ta{அகழி} mean? (\textquotedblleft{}A moat\textquotedblright{}.) Say it once more.
\end{tcolorbox}
\section[\texorpdfstring{\ta{சுரங்கம்} (curaṅkam)}{சுரங்கம் (curaṅkam)}]{\ta{சுரங்கம்} (curaṅkam) — a tunnel; a mine}

\label{lesson:TA-C335-curankam}



\begin{quote}
Before the new one: say the Tamil for to besiege, then the Tamil for a moat.
\end{quote}

\subsection*{You'll want to know: \ta{சுரங்கம்}}
\textbf{\ta{சுரங்கம்}} — \emph{curaṅkam} — \textquotedblleft{}a tunnel; a mine\textquotedblright{}.

\begin{quote}
\ta{வீரர்கள்} \ta{சுவருக்கு} \ta{அடியில்} \ta{ஒரு} \ta{சுரங்கம்} \ta{தோண்டினார்கள்}.
\end{quote}

\begin{quote}
\emph{vīrarkaḷ cuvarukku aṭiyil oru curaṅkam tōṇṭiṉārkaḷ.}
\end{quote}

\begin{quote}
The soldiers dug a tunnel under the wall.
\end{quote}

Also a mine: \ta{நிலக்கரிச்} \ta{சுரங்கம்}, a coal mine.

\begin{cousinweb}
\textbf{\ta{சுரங்கம்}} (\emph{curaṅkam}) is borrowed from Sanskrit \emph{suruṅgā}, \textquotedblleft{}a hole made underground\textquotedblright{}, for war (Monier-Williams).
\end{cousinweb}

\begin{grammarlens}[title={What you've built}]
One more word for telling what happened.
\end{grammarlens}

\subsection*{Your turn}
\begin{itemize}
\raggedright
  \item \emph{Say it:} \emph{curaṅkam}
  \item \emph{Say it:} \emph{curaṅkam}, once more
  \item \emph{Say it:} \emph{vīrarkaḷ cuvarukku aṭiyil oru curaṅkam tōṇṭiṉārkaḷ.} — The soldiers dug a tunnel under the wall.
  \item \emph{Recall:} say the Tamil for to besiege, then the Tamil for a moat, then say \emph{curaṅkam} again
\end{itemize}

\begin{tcolorbox}[breakable,colback=teal!4,colframe=teal!35!black,title={Before you move on}]
What does \ta{சுரங்கம்} mean? (\textquotedblleft{}A tunnel; a mine\textquotedblright{}.) Say it once more.
\end{tcolorbox}
\section[\texorpdfstring{\ta{அரண்} (araṇ)}{அரண் (araṇ)}]{\ta{அரண்} (araṇ) — a fortification; a defence}

\label{lesson:TA-C335-aran}



\begin{quote}
Before the new one: say the Tamil for a moat, then the Tamil for a tunnel.
\end{quote}

\subsection*{You'll want to know: \ta{அரண்}}
\textbf{\ta{அரண்}} — \emph{araṇ} — \textquotedblleft{}a fortification; a defence\textquotedblright{}.

\begin{quote}
\ta{அந்த} \ta{மலை} \ta{நகரத்துக்கு} \ta{ஒரு} \ta{இயற்கை} \ta{அரணாக} \ta{இருந்தது}.
\end{quote}

\begin{quote}
\emph{anta malai nakarattukku oru iyaṟkai araṇāka iruntatu.}
\end{quote}

\begin{quote}
That hill was a natural defence for the city.
\end{quote}

-\ta{ஆக}, as: \ta{அரணாக} \ta{இருந்தது}, it served as a defence.

\begin{cousinweb}
\textbf{\ta{அரண்}} (\emph{araṇ}) is usually traced to Sanskrit \emph{śaraṇa}, \textquotedblleft{}protecting, shelter\textquotedblright{} (CDIAL 12326), as DEDR 201 notes.
\end{cousinweb}

\begin{grammarlens}[title={What you've built}]
One more word for telling what happened.
\end{grammarlens}

\subsection*{Your turn}
\begin{itemize}
\raggedright
  \item \emph{Say it:} \emph{araṇ}
  \item \emph{Say it:} \emph{araṇ}, once more
  \item \emph{Say it:} \emph{anta malai nakarattukku oru iyaṟkai araṇāka iruntatu.} — That hill was a natural defence for the city.
  \item \emph{Recall:} say the Tamil for a moat, then the Tamil for a tunnel, then say \emph{araṇ} again
\end{itemize}

\begin{tcolorbox}[breakable,colback=teal!4,colframe=teal!35!black,title={Before you move on}]
What does \ta{அரண்} mean? (\textquotedblleft{}A fortification; a defence\textquotedblright{}.) Say it once more.
\end{tcolorbox}
\section[\texorpdfstring{\ta{ஒற்றன்} (oṟṟaṉ)}{ஒற்றன் (oṟṟaṉ)}]{\ta{ஒற்றன்} (oṟṟaṉ) — a spy}

\label{lesson:TA-C335-orran}



\begin{quote}
Before the new one: say the Tamil for a tunnel, then the Tamil for a fortification.
\end{quote}

\subsection*{You'll want to know: \ta{ஒற்றன்}}
\textbf{\ta{ஒற்றன்}} — \emph{oṟṟaṉ} — \textquotedblleft{}a spy\textquotedblright{}.

\begin{quote}
\ta{ஒரு} \ta{ஒற்றன்} \ta{மன்னருக்கு} \ta{ரகசியச்} \ta{செய்தி} \ta{கொண்டுவந்தான்}.
\end{quote}

\begin{quote}
\emph{oru oṟṟaṉ maṉṉarukku rakaciyac ceyti koṇṭuvantāṉ.}
\end{quote}

\begin{quote}
A spy brought the king a secret message.
\end{quote}

From \ta{ஒற்று}, to spy out.

\begin{cousinweb}
\textbf{\ta{ஒற்றன்}} (\emph{oṟṟaṉ}) is native Dravidian: DEDR 1022 gives it under \ta{ஒற்று}, \textquotedblleft{}to spy out\textquotedblright{}.
\end{cousinweb}

\begin{grammarlens}[title={What you've built}]
One more word for telling what happened.
\end{grammarlens}

\subsection*{Your turn}
\begin{itemize}
\raggedright
  \item \emph{Say it:} \emph{oṟṟaṉ}
  \item \emph{Say it:} \emph{oṟṟaṉ}, once more
  \item \emph{Say it:} \emph{oru oṟṟaṉ maṉṉarukku rakaciyac ceyti koṇṭuvantāṉ.} — A spy brought the king a secret message.
  \item \emph{Recall:} say the Tamil for a tunnel, then the Tamil for a fortification, then say \emph{oṟṟaṉ} again
\end{itemize}

\begin{tcolorbox}[breakable,colback=teal!4,colframe=teal!35!black,title={Before you move on}]
What does \ta{ஒற்றன்} mean? (\textquotedblleft{}A spy\textquotedblright{}.) Say it once more.
\end{tcolorbox}
